% Zone „Das kennst du schon“ – aus bank/terme/zone.jsonl (zusammenbau v0.1)
\einheitenkopf[zone][Kennst du schon]{Das kennst du schon}
\setcounter{aufgabe}{0}
%% TODO Grundvorstellungs-Aufgabe als erste Hauptnummer der Zone (2.8) fehlt in der Bank

%% TODO Titel: Ich-kann-Satz fehlt in der Bank (Platzhalter: Kettenname) – Nr. 1
%% TODO Anweisung: ein Satz über der ersten Teilaufgabe fehlt in der Bank – Nr. 1
\begin{aufgabe}{Addieren und Subtrahieren negativer Zahlen}
%% TODO Darstellung für schwach fehlt in der Bank (terme-zone-f1-v1; Abschnitt „Für schwache Schüler“)
\swz{Rechne: $4 - 9$}{}
%% TODO Darstellung für schwach fehlt in der Bank (terme-zone-f1-v3; Abschnitt „Für schwache Schüler“)
\swz{Rechne: $-2{,}5 + 1$}{}
\end{aufgabe}

%% TODO Titel: Ich-kann-Satz fehlt in der Bank (Platzhalter: Kettenname) – Nr. 2
%% TODO Anweisung: ein Satz über der ersten Teilaufgabe fehlt in der Bank – Nr. 2
\begin{aufgabe}{Multiplizieren mit Vorzeichen}
%% TODO Darstellung für schwach fehlt in der Bank (terme-zone-f2-v1; Abschnitt „Für schwache Schüler“)
\swz{Rechne: $(-5) \cdot 2$}{}
%% TODO Darstellung für schwach fehlt in der Bank (terme-zone-f2-v3; Abschnitt „Für schwache Schüler“)
\swz{Rechne: $(-0{,}5) \cdot 8$}{}
\end{aufgabe}

%% TODO Titel: Ich-kann-Satz fehlt in der Bank (Platzhalter: Kettenname) – Nr. 3
%% TODO Anweisung: ein Satz über der ersten Teilaufgabe fehlt in der Bank – Nr. 3
\begin{aufgabe}{Dezimalzahlen und einfache Brüche als Vorzahlen}
%% TODO Darstellung für schwach fehlt in der Bank (terme-zone-f3-v1; Abschnitt „Für schwache Schüler“)
\swz{Rechne: $1{,}5 + 2{,}5$}{}
%% TODO Darstellung für schwach fehlt in der Bank (terme-zone-f3-v3; Abschnitt „Für schwache Schüler“)
\swz{Rechne: $0{,}5 - 1{,}5$}{}
\end{aufgabe}

%% TODO Titel: Ich-kann-Satz fehlt in der Bank (Platzhalter: Kettenname) – Nr. 4
%% TODO Anweisung: ein Satz über der ersten Teilaufgabe fehlt in der Bank – Nr. 4
\begin{aufgabe}{Punkt vor Strich mit Zahlen}
%% TODO Darstellung für schwach fehlt in der Bank (terme-zone-f4-v1; Abschnitt „Für schwache Schüler“)
\swz{Rechne: $5 + 2 \cdot 4$}{}
%% TODO Darstellung für schwach fehlt in der Bank (terme-zone-f4-v3; Abschnitt „Für schwache Schüler“)
\swz{Rechne: $6 - 4 \cdot 2{,}5$}{}
\end{aufgabe}

%% TODO Titel: Ich-kann-Satz fehlt in der Bank (Platzhalter: Kettenname) – Nr. 5
%% TODO Anweisung: ein Satz über der ersten Teilaufgabe fehlt in der Bank – Nr. 5
\begin{aufgabe}{Addieren und Subtrahieren negativer Zahlen – Fehler finden}
\swz[3]{Finde den Fehler und rechne richtig: \rechnung{-7 - 5 &= -2}}{}
\end{aufgabe}

%% TODO Titel: Ich-kann-Satz fehlt in der Bank (Platzhalter: Kettenname) – Nr. 6
%% TODO Anweisung: ein Satz über der ersten Teilaufgabe fehlt in der Bank – Nr. 6
\begin{aufgabe}{Addieren und Subtrahieren negativer Zahlen – selbst rechnen}
%% TODO Darstellung für schwach fehlt in der Bank (terme-zone-f1-v6; Abschnitt „Für schwache Schüler“)
\swz{Rechne: $-8 - 6$}{}
\end{aufgabe}
